%% ma: L10-B14-0071
%% lop: 10
%% chuong: 5
%% bai: 14
%% dang: 10.14.7
%% loai: TN4
%% muc: VD
%% dap_an: "D"
%% nguon: "chuyen-de-cac-so-dac-trung-cua-mau-so-lieu-khong-ghep-nhom-toan-10-chuong-trinh-moi.pdf (GV. Phan Nhật Linh) (P05), Bài 03 Dạng 2 Phần I Câu 15"
%% kien_thuc: "Lớp 10 Bài 13 (số trung bình), Bài 14 (phương sai, độ lệch chuẩn)"
%% trang_thai: chinh-thuc
%% ghi_chu: "Tự duyệt theo yêu cầu giáo viên 22/10/2026."
\begin{ex}
Lương (đơn vị: triệu đồng) trong tháng 2 của các công nhân trong một tổ sản xuất như sau:
\begin{center}
\begin{tabular}{|c|c|c|c|c|c|c|c|c|c|}\hline
$5$ & $15$ & $5$ & $10$ & $6$ & $6$ & $7$ & $15$ & $7$ & $6$\\\hline
$6$ & $7$ & $7$ & $15$ & $10$ & $7$ & $6$ & $10$ & $7$ & $10$\\\hline
\end{tabular}
\end{center}
Độ lệch chuẩn của mẫu số liệu gần bằng
\choice{$3{,}50$}{$3$}{$10{,}23$}{\True $3{,}20$}
\loigiai{Bảng tần số: giá trị $5$; $6$; $7$; $10$; $15$ có tần số $2$; $5$; $6$; $4$; $3$ ($n=20$).

$\overline x=\dfrac{167}{20}=8{,}35$; $s^2=10{,}2275$; $s\approx3{,}20$. Chọn D.}
\end{ex}
